\section{Kiểm thử hệ thống}
\subsection{Kiểm thử Web}
Môi trường kiểm thử:
\begin{enumerate}[-]
    \item Hệ điều hành: Windows
    \item Phiên bản hệ điều hành: 23H2 (OS Build 22631.3593)
    \item Phiên bản trình duyệt Chrome: 124.0.6367.208
\end{enumerate}
Đầu tiên là bảng Kiểm thử các yêu cầu chức năng của trang Dashboard, nơi hiển thị tổng số đơn hàng theo từng trạng thái, tổng giá trị đơn hàng theo từng trạng thái, doanh thu và các sản phẩm bán chạy.
\begin{longtable}{|p{4cm}|p{4cm}|p{4cm}|p{2cm}|}
    \hline
    Chức năng & Tiền điều kiện & Kết quả dự kiến & Kết quả \\
    \hline
    A1. Hiển thị tổng số đơn hàng theo từng trạng thái & Truy cập vào trang Dashboard & Hiển thị tổng số đơn hàng theo từng trạng thái & Đạt\\
    \hline
     A2. Hiển thị tổng giá trị đơn hàng theo từng trạng thái & Truy cập vào trang Dashboard & Hiển thị tổng giá trị đơn hàng theo từng trạng thái & Đạt\\
    \hline
     A3. Hiển thị doanh thu tháng này & Truy cập vào trang Dashboard, ở phần doanh thu có chế độ mặc định "Tháng này"  & Hiển thị doanh thu tháng này & Đạt\\
     \hline
     A4. Hiển thị doanh thu tháng trước & Truy cập vào trang Dashboard, ở phần doanh thu chọn chế độ "Tháng trước" & Hiển thị doanh thu tháng trước & Đạt \\
     \hline
     A5. Hiển thị doanh thu trong một khoảng thời gian & Truy cập vào trang Dashboard, ở phần doanh thu chọn chế độ "Tùy chỉnh khoảng thời gian" và chọn khoảng thời gian & Hiển thị doanh thu trong khoảng thời gian đã chọn & Đạt \\
     \hline
     A6. Hiển thị danh sách các sản phẩm bán chạy trong tháng này & Truy cập vào trang Dashboard, ở phần sản phẩm bán chạy có chế độ mặc định là "Tháng này" & Hiển thị danh sách các sản phẩm bán chạy trong tháng này & Đạt \\
     \hline 
     A7. Hiển thị danh sách các sản phẩm bán chạy trong tháng trước & Truy cập vào trang Dashboard, ở phần sản phẩm bán chạy chọn chế độ "Tháng này" & Hiển thị danh sách các sản phẩm bán chạy trong tháng này & Đạt \\
     \hline 
     A8. Hiển thị danh sách các sản phẩm bán chạy trong tuần này & Truy cập vào trang Dashboard, ở phần sản phẩm bán chạy chọn chế độ "Tuần này" & Hiển thị danh sách các sản phẩm bán chạy trong tuần này & Đạt \\
     \hline 
     A9. Hiển thị danh sách các sản phẩm bán chạy trong tuần trước & Truy cập vào trang Dashboard, ở phần sản phẩm bán chạy chọn chế độ "Tuần trước" & Hiển thị danh sách các sản phẩm bán chạy trong tuần trước & Đạt \\
    \hline
    \caption{Bảng kiểm thử yêu cầu chức năng trang Dashboard}
\end{longtable}
Tiếp theo là bảng Kiểm thử các yêu cầu chức năng của trang thể loại, bao gồm hiển thị các thể loại theo dạng cây thư mục, chức năng tạo, chỉnh sửa, xóa thể loại.
\begin{longtable}{|p{4cm}|p{4cm}|p{4cm}|p{2cm}|}
    \hline
    Chức năng & Tiền điều kiện & Kết quả dự kiến & Kết quả \\
    \hline
    B1. Hiển thị danh sách tất cả các thể loại& Truy cập vào trang thể loại & Hiển thị danh sách tất cả các thể loại theo dạng cây thư mục & Đạt\\
    \hline
     B2. Tạo mới thể loại & Nhấn vào button thêm thể loại và nhập các thông tin cho thể loại mới  & Hiển thị thể loại mới trong cây thư mục đúng vị trí & Đạt\\
    \hline
     B3. Cập nhật thể loại & Nhấn vào icon Edit và chỉnh sửa các thông tin của thể loại & Hiển thị lại danh sách thể loại theo đúng cấu trúc cây thư mục sau khi vừa cập nhật thể loại  & Đạt\\
    \hline
    B4. Xóa thể loại & Nhấn vào icon Delete và xác nhận có muốn xóa thể loại vừa nhấn hay không & Hiển thị lại danh sách thể loại theo đúng cấu trúc cây thư mục sau khi vừa xóa thể loại & Đạt\\
    \hline
    \caption{Bảng kiểm thử các yêu cầu chức năng của trang thể loại}
\end{longtable}
Tiếp theo là bảng Kiểm thử các yêu cầu chức năng của trang sản phẩm, nơi hiển thị danh sách tất cả các sản phẩm, hỗ trợ tính năng lọc/tìm kiếm các sản phẩm, hiển thị chi tiết các sản phẩm, tạo mới/cập nhật/ sản phẩm
\begin{longtable}{|p{4cm}|p{4cm}|p{4cm}|p{2cm}|}
    \hline
    Chức năng & Tiền điều kiện & Kết quả dự kiến & Kết quả \\
    \hline
    C1. Hiển thị danh sách tất cả các sản phẩm  & Truy cập vào trang sản phẩm & Hiển thị danh sách tất cả các sản phẩm & Đạt\\
    \hline
     C2. Hiển thị chi tiết sản phẩm  & Truy cập vào trang sản phẩm và nhấn xem chi tiết một sản phẩm & Hiển thị các thông tin chi tiết của sản phẩm & Đạt\\
    \hline
     C3. Tạo mới sản phẩm & Nhấn vào button thêm sản phẩm và nhập các thông tin cho sản phẩm mới & Tạo mới sản phẩm thành công và hiển thị form cho phép có thể thêm variant cho sản phẩm mới & Đạt\\
    \hline
    C4. Cập nhật sản phẩm & Nhấn button cập nhật ở trang chi tiết sản phẩm và chỉnh sửa các thông tin sản phẩm & Hiển thị lại thông tin sản phẩm sau khi đã cập nhật & Đạt\\
    \hline
    C5. Tìm kiếm sản phẩm & Nhập từ khóa vào ô tìm kiếm & Hiển thị các sản phẩm có tên chứa từ khóa đã nhập & Đạt \\
    \hline
    C6. Lọc sản phẩm theo min price và max price & Nhập phạm vi giá cả cho sản phẩm & Hiển thị các sản phẩm nằm trong phạm vi min price và max price & Đạt  \\
    \hline
    C7. Lọc sản phẩm theo số sao đánh giá sản phẩm & Chọn số sao đánh giá cho sản phẩm & Hiển thị các sản phẩm tương ứng với số sao đã chọn  & Đạt\\
    \hline
    C8. Lọc sản phẩm theo thể loại & Chọn thể loại muốn lọc & Hiển thị các sản phẩm tương ứng với thể loại đã chọn & Đạt \\
    \hline
    C9. Sắp xếp sản phẩm & Chọn chế độ muốn sắp xếp sản phẩm & Hiển thị các sản phẩm theo chế độ sắp xếp & Đạt\\
    \hline
    \caption{Bảng kiểm thử các yêu cầu chức năng của trang sản phẩm}
\end{longtable}
Tiếp theo là bảng Kiểm thử các yêu cầu chức năng của Review, bao gồm xem tất cả đánh giá của từng sản phẩm và xóa đánh giá sản phẩm.
\begin{longtable}{|p{4cm}|p{4cm}|p{4cm}|p{2cm}|}
    \hline
    Chức năng & Tiền điều kiện & Kết quả dự kiến & Kết quả \\
    \hline
    D1. Hiển thị tất cả các đánh giá của một sản phẩm  & Xem chi tiết một sản phẩm và chọn vào tab đánh giá sản phẩm & Hiển thị tất cả các đánh giá của sản phẩm đó & Đạt\\
    \hline
     D2. Xóa đánh giá sản phẩm  & Nhấn icon Delete tại đánh giá sản phẩm muốn xóa & Danh sách đánh giá sản phẩm được cập nhật lại sau khi vừa xóa đánh giá & Đạt\\
    \hline
    \caption{Bảng kiểm thử các yêu cầu chức năng của trang variant}
\end{longtable}
Tiếp theo là bảng Kiểm thử các yêu cầu chức năng của trang Variant, nơi hiển thị danh sách tất cả các variant theo từng sản phẩm, tạo mới/chỉnh sửa/cập nhật/xóa các variant.
\begin{longtable}{|p{4cm}|p{4cm}|p{4cm}|p{2cm}|}
    \hline
    Chức năng & Tiền điều kiện & Kết quả dự kiến & Kết quả \\
    \hline
    E1. Hiển thị danh sách các variant theo sản phẩm  & Truy cập vào trang sản phẩm và nhấn xem chi tiết một sản phẩm & Hiển thị danh sách variant theo sản phẩm dưới dạng bảng & Đạt\\
    \hline
     E2. Thêm variant cho sản phẩm & Nhấn button thêm variant và nhập các thông tin cho variant mới & Variant mới được hiển thị trong bảng variant của sản phẩm  & Đạt\\
    \hline
     E3. Cập nhật variant của sản phẩm& Nhấn icon Edit của variant muốn chỉnh sửa & Bảng variant cập nhật thông tin cho variant vừa được chỉnh sửa  & Đạt\\
     E4. Xóa variant của sản phẩm & Nhấn icon Delete của variant muốn xóa & Bảng variant không còn chứa variant vừa xóa & Đạt \\
    \hline
    \caption{Bảng kiểm thử các yêu cầu chức năng của trang variant}
\end{longtable}
Tiếp theo là bảng Kiểm thử các yêu cầu chức năng của trang mã giảm giá, nơi hiển thị danh sách tất cả các mã giảm giá, tạo mới/chỉnh sửa/cập nhật/xóa các mã giảm giá, gán mã giảm giá cho thể loại/sản phẩm/variant/người dùng.
\begin{longtable}{|p{4cm}|p{4cm}|p{4cm}|p{2cm}|}
    \hline
    Chức năng & Tiền điều kiện & Kết quả dự kiến & Kết quả \\
    \hline
    F1. Hiển thị danh sách tất cả mã giảm giá & Truy cập vào trang mã giảm giá & Hiển thị danh sách tất cả mã giảm giá dưới dạng bảng & Đạt\\
    \hline
     F2. Tạo mới mã giảm giá & Nhấn vào button thêm mã giảm giá trong trang mã giảm giá và nhập các thông tin cho mã giảm giá mới & Mã giảm giá mới được hiển thị trong bảng & Đạt\\
    \hline
     F3. Hiển thị thông tin chi tiết của mã giảm giá & Nhấn vào button Xem chi tiết của một mã giảm giá trong bảng  & Hiển thị thông tin chi tiết của mã giảm giá & Đạt\\
     \hline
     F4. Chỉnh sửa thông tin của mã giảm giá & Nhấn vào icon Edit của một mã giảm giá trong bảng và chỉnh sửa các thông tin & Hiển thị thông tin sau khi cập nhật của mã giảm giá trong bảng & Đạt\\
     \hline
     F5. Xóa mã giảm giá & Nhấn vào button Delete của một mã giảm giá & Mã giảm giá xóa thành công khỏi bảng & Đạt \\
     \hline
     F6. Cập nhật thể loại cho mã giảm giá & Nhấn vào button cập nhật thể loại và chọn thể loại tương ứng & Mã giảm giá đã cập nhật thể loại & Đạt\\
     \hline
     F7. Cập nhật người dùng cho mã giảm giá & Nhấn vào button Gán mã giảm giá cho người dùng và chọn các người dùng được áp dụng & Mã giảm giá đã được áp dụng cho các người dùng được chọn & Đạt \\
    \hline
    \caption{Bảng kiểm thử các yêu cầu chức năng của trang mã giảm giá}
\end{longtable}
Tiếp theo là bảng Kiểm thử các yêu cầu chức năng của trang đơn hàng, nơi hiển thị danh sách tất cả các đơn hàng theo trạng thái, xem chi tiết từng đơn hàng, cập nhật trạng thái đơn hàng.
\begin{longtable}{|p{4cm}|p{4cm}|p{4cm}|p{2cm}|}
    \hline
    Chức năng & Tiền điều kiện & Kết quả dự kiến & Kết quả \\
    \hline
    G1. Hiển thị danh sách tất cả các đơn hàng & Truy cập vào trang đơn hàng & Hiển thị danh sách tất cả các đơn hàng theo trạng thái & Đạt\\
    \hline
     G2. Hiển thị thông tin chi tiết của đơn hàng & Truy cập vào trang đơn hàng và nhấn button Xem chi tiết vào một đơn hàng & Hiển thị thông tin chi tiết đơn hàng đó & Đạt\\
    \hline
     G3. Hiển thị các sản phẩm trong đơn hàng & Truy cập vào trang đơn hàng và nhấn button Xem chi tiết vào một đơn hàng  & Hiển thị các sản phẩm trong đơn hàng & Đạt\\
     \hline
     G4. Hiển thị giá trị của đơn hàng & Truy cập vào trang đơn hàng và nhấn button Xem chi tiết vào một đơn hàng  & Hiển thị giá trị của đơn hàng & Đạt\\
      \hline
     G5. Cập nhật trạng thái của đơn hàng & Truy cập vào trang đơn hàng và nhấn button Xem chi tiết vào một đơn hàng, chọn trạng thái phù hợp để cập nhật  & Cập nhật trạng thái của đơn hàng & Đạt\\
    \hline
    \caption{Bảng kiểm thử các yêu cầu chức năng của trang đơn hàng}
\end{longtable}
Tiếp theo là bảng Kiểm thử các yêu cầu chức năng của trang khách hàng, nơi hiển thị danh sách tất cả các khách hàng, xem thông tin cá nhân, thông tin nhận hàng, các đơn hàng đã đặt, các đánh giá sản phẩm của từng khách hàng.
\begin{longtable}{|p{4cm}|p{4cm}|p{4cm}|p{2cm}|}
    \hline
    Chức năng & Tiền điều kiện & Kết quả dự kiến & Kết quả \\
    \hline
    H1. Hiển thị danh sách tất cả các khách hàng & Truy cập vào trang khách hàng & Hiển thị danh sách tất cả các khách hàng & Đạt\\
    \hline
     H2. Hiển thị thông tin cá nhân của khách hàng & Nhấn vào button Xem chi tiết của một khách hàng trong bảng & Hiển thị thông tin cá nhân của một khách hàng & Đạt\\
    \hline
     H3. Hiển thị thông tin nhận hàng của khách hàng & Chuyển qua tab "Thông tin nhận hàng" của trang xem thông tin cá nhân của khách hàng & Hiển thị thông tin nhận hàng của khách hàng & Đạt\\
     \hline
     H4. Hiển thị các đơn hàng của một khách hàng & Chuyển qua tab "Thông tin nhận hàng" của trang xem thông tin cá nhân của khách hàng & Hiển thị các đơn hàng của khách hàng & Đạt\\
     \hline
      H5. Hiển thị các đánh giá sản phẩm của một khách hàng & Chuyển qua tab "Đánh giá sản phẩm" của trang xem thông tin cá nhân của khách hàng & Hiển thị các đánh giá sản phẩm của khách hàng & Đạt\\
    \hline
    \caption{Bảng kiểm thử các yêu cầu chức năng của trang khách hàng}
\end{longtable}
Tiếp theo là bảng Kiểm thử các yêu cầu chức năng của trang nhân viên, nơi hiển thị danh sách tất cả các nhân viên, xem thông tin cá nhân của từng nhân viên.
\begin{longtable}{|p{4cm}|p{4cm}|p{4cm}|p{2cm}|}
    \hline
    Chức năng & Tiền điều kiện & Kết quả dự kiến & Kết quả \\
    \hline
    I1. Hiển thị danh sách tất cả các nhân viên & Truy cập vào trang nhân viên & Hiển thị danh sách tất cả các nhân viên & Đạt\\
    \hline
     I2. Hiển thị thông tin chi tiết của một nhân viên & Nhấn vào button Xem chi tiết của một nhân viên trong bảng  & Hiển thị thông tin chi tiết của nhân viên & Đạt\\
    \hline
    \caption{Bảng kiểm thử các yêu cầu chức năng của trang nhân viên}
\end{longtable}
\newpage
\subsection{Kiểm thử Ứng dụng di động}
Môi trường kiểm thử:
\begin{enumerate}[-]
    \item Hệ điều hành: Android
    \item Phiên bản hệ điều hành: Android 14
\end{enumerate}
\noindent Chức năng đăng nhập/đăng ký
\begin{longtable}{|p{4cm}|p{4cm}|p{4cm}|p{2cm}|}
    \hline
    Chức năng & Tiền điều kiện & Kết quả dự kiến & Kết quả \\
    \hline
    A1. Đăng nhập bằng tài khoản, mật khẩu đúng & Nhập tên đăng nhập và mật khẩu & Hiển thị thông báo đăng nhập thành công, chuyển đến trang chủ & Đạt\\
    \hline
    A2. Đăng nhập bằng tài khoản không hợp lệ & Nhập tên đăng nhập và mật khẩu không hợp lệ & Hiển thị thông báo không tìm thấy tài khoản & Đạt\\
    \hline
    A3. Đăng nhập bỏ trống tài khoản & Nhập mật khẩu & Vô hiệu hoá nút Đăng nhập & Đạt\\
    \hline
    A4. Đăng nhập bỏ trống mật khẩu & Nhập tài khoản & Vô hiệu hoá nút Đăng nhập & Đạt\\
    \hline
    A5. Đăng nhập bằng tài khoản đúng nhưng mật khẩu sai & Nhập tài khoản, và mật khẩu không đúng & Hiển thị thông báo nhập sai mật khẩu & Đạt\\
    \hline
    A6. Đăng nhập bằng Google & Chọn tuỳ chọn đăng nhập bàng Google & Hiển thị thông báo đăng nhập thành công với tài khoản Google & Đạt\\
    \hline
     A7. Đăng ký bằng tài khoản, mật khẩu & Nhập mật khẩu, số điện thoại và email chưa được sử dụng & Hiển thị thông báo đăng ký thành công & Đạt\\
    \hline
    A8. Đăng ký bằng số điện thoại đã tồn tại & Nhập mật khẩu, số điện thoại đã tồn tại và email chưa được sử dụng & Hiển thị thông báo số điện thoại đã được đăng ký & Đạt\\
    \hline
    A9. Đăng ký bằng email đã tồn tại & Nhập mật khẩu, số điện thoại chưa sử dụng và email đã tồn tại & Hiển thị thông báo email đã được đăng ký & Đạt\\
    \hline
    A10. Đăng ký bằng email và số điện thoại đã tồn tại & Nhập mật khẩu, số điện thoại chưa sử dụng và email đã tồn tại & Hiển thị thông báo email đã được đăng ký & Đạt\\
    \hline
    A11. Đăng ký bằng tài khoản, mật khẩu bỏ trống các trường thông tin & Không nhập một số trường thông tin & Vô hiệu hoá nút đăng ký & Đạt\\
    \hline
    \caption{Kiểm thử yêu cầu chức năng trang đăng nhập và đăng ký}
\end{longtable}

\noindent Chức năng tìm kiếm, lọc sản phẩm, nơi người dùng có thể tìm kiếm sản phẩm theo tên, hình ảnh, sắp xếp sản phẩm. Người dùng cần phải bấm vào ô tìm kiếm hoặc icon hình camera bên cạnh ô tìm kiếm để thực hiện chức năng.

\begin{longtable}{|p{4cm}|p{4cm}|p{4cm}|p{2cm}|}
    \hline
    Chức năng & Tiền điều kiện & Kết quả dự kiến & Kết quả \\
    \hline
    B1. Tìm kiếm sản phẩm theo tên có tồn tại & Nhập tên một sản phẩm tồn tại  & Hiển thị đúng sản phẩm theo tên & Đạt\\
    \hline
    B2. Tìm kiếm sản phẩm bằng hình ảnh & Nhấn vào biểu tượng Camera, tài lên 1 hình ảnh & Hiển thị những sản phẩm được tìm thấy bằng hình ảnh & Đạt\\
    \hline
    B3. Tìm kiếm sản phẩm theo danh mục & Dưới thanh tìm kiếm hiển thị một số danh mục, và chọn vào 1 danh mục & Hiển thị những sản phẩm theo danh mục đã chọn & Đạt\\
    \hline
    B4. Sắp xếp những sản phẩm theo giá & Chọn vào mục Sắp xếp, sau đó chọn giắ tăng dần hoặc giảm dần & Hiển thị những sản phẩm được sắp xếp theo giá tăng dần hoặc giảm dần & Đạt\\
    \hline
    B5. Sắp xếp những sản phẩm theo phổ biến & Chọn vào mục Sắp xếp, sau đó phổ biến nhất & Hiển thị những sản phẩm được sắp xếp theo độ phổ biến & Đạt\\
    \hline
    B6. Sắp xếp những sản phẩm mới nhất & Chọn vào mục Sắp xếp, sau đó chọn Mới nhất & Hiển thị những sản phẩm được sắp xếp theo mới nhất & Đạt\\
    \hline
     B7. Sắp xếp những sản phẩm được giảm giá nhiều nhất & Chọn vào mục Sắp xếp, sau đó chọn Giảm giá nhiều nhất & Hiển thị những sản phẩm được sắp xếp theo mức giảm giá nhiều nhất & Đạt\\
    \hline
    B8. Lọc sản phẩm theo giá tiền & Điền giá tiền nhỏ nhất và lớn nhất & Hiển thị những sản phẩm nằm trong giá tiền & Đạt\\
    \hline
    B9. Lọc sản phẩm theo số sao đánh giá & Chọn mức sao sản phẩm & Hiển thị những sản phẩm có số sao lớn hơn hoặc bằng mức đã chọn & Đạt\\
    \hline
    B10. Lọc sản phẩm theo mã giảm giá hiện có & Chọn mã voucher cần lọc & Hiển thị những sản phẩm có thể áp dụng được voucher & Đạt\\
    \hline
    \caption{Kiểm thử yêu cầu chức năng tìm kiếm sản phẩm}
\end{longtable}

\noindent Chức năng quản lý giỏ hàng, bao gồm việc thêm sản phẩm vào giỏ hàng hoặc xoá, để có thể thêm sản phẩm vào được giỏ hàng thì bắt buộc người dùng phải đăng nhập trước.

\begin{longtable}{|p{4cm}|p{4cm}|p{4cm}|p{2cm}|}
    \hline
    Chức năng & Tiền điều kiện & Kết quả dự kiến & Kết quả \\
    \hline
    C1. Thêm một sản phẩm vào giỏ hàng đã chọn kích thước, màu sắc  & Chọn kích thước, màu sắc và chọn thêm vào giỏ & Hiển thị thông báo Thêm vào giỏ hàng thành công  & Đạt\\
    \hline
    C2. Thêm một sản phẩm vào giỏ hàng mà không chọn kích thước hoặc màu sắc & Không chọn kích thước hoặc màu sắc & Vô hiệu hoá nút thêm vào giỏ hàng & Đạt\\
    \hline
    C3. Thêm một sản phẩm mà số lượng đã hết & Chọn vào sản phẩm có số lượng là 0 & Vô hiệu hoá việc chọn sản phẩm có số lượng bằng 0 và Thêm vào giỏ & Đạt\\
    \hline
    C4. Thêm sản phẩm vào giỏ hàng mà chưa đăng nhập & Chọn sản phẩm, kích thước và màu sắc, sau đó chọn thêm vào giỏ hàng & Hiển thị thông báo yêu cầu đăng nhập & Đạt \\
    \hline
    C5. Xoá sản phẩm ra khỏi giỏ hàng & Ở trang giỏ hàng, người dùng bấm xoá sản phẩm & Sản phẩm được xoá khỏi giỏ hàng  & Đạt \\
    \hline
    C6. Xem sản phẩm có trong giỏ hàng & Chọn vào giỏ hàng & Hiển thị những sản phảm trong giỏ hàng & Đạt \\
    \hline
    C7. Mua hàng mà không chọn sản phẩm nào trong giỏ hàng & Bỏ chọn tất cả sản phẩm trong giỏ hàng & Vô hiệu hoá nút mua hàng & Đạt \\
    \hline
    C8. Mua hàng mà đã chọn một trong những sản phẩm trong giỏ hàng & Chọn sản phẩm trong giỏ hàng và bấm mua hàng & Chuyển đến trang đơn hàng & Đạt \\
    \hline
    \caption{Kiểm thử yêu cầu chức năng giỏ hàng}
\end{longtable}

\noindent Chức năng về sản phẩm yêu thích. Nơi sẽ lưu những sản phầm mà người dùng mong muốn sẽ mua trong tương lai.
\begin{longtable}{|p{4cm}|p{4cm}|p{4cm}|p{2cm}|}
    \hline
    Chức năng & Tiền điều kiện & Kết quả dự kiến & Kết quả \\
    \hline
    D1. Xem sản phẩm yêu thích & Chọn vào trang yêu thích & Hiên thị những sản phẩm đã đánh dấu yêu thích & Đạt \\
    \hline
    D2. Lọc sản phẩm theo danh mục & Ở trang yêu thích, chọn lọc theo danh mục & Sản phẩm yêu thích được lọc theo danh mục. & Đạt \\
    \hline
    D3. Thêm sản phẩm yêu thích & Bấm vào icon trái tim  & Hiện thông báo sản phẩm đã được thêm vào yêu thích & Đạt \\
    \hline
    D4. Xoá sản phẩm khỏi yêu thích & Bấm vào icon trái tim & Hiện thông bảo sản phẩm đã đượpc xoá khỏi yêu thích & Đạt \\
    \hline
    \caption{Kiểm thử yêu cầu chức năng sản phẩm yêu thích}
\end{longtable}



\newpage
%%%%%%%%%%%%%%%%%%%%%%%%%%%%%%%%%%%%%%%%%%%
\subsection{Kiểm thử API Backend}


\noindent Postman là một công cụ kiểm thử API được nhiều người ưa chuộng vì tính đa năng và dễ sử dụng. Sau khi triển khai server, chúng tôi đã tiến hành kiểm thử lên server có cấu hình như sau:
\begin{itemize}
    \item 1 CPU core
    \item 2 GB memory
    \item 25 GB SSD
    \item 1000 GB transfer
\end{itemize}
Tiếp tục, chúng tôi sử dụng performance test của postman để test các collection. Với cấu hình performance test như sau:
\begin{itemize}
    \item Virtual users: 100
    \item Test duration: 5 mins
    \item Load profile: Fixed
\end{itemize}
Kết quả chạy của các phiên như sau:

\begin{table}[H]
\centering
\begin{tabularx}{\textwidth}{|>{\centering\arraybackslash}X|>{\centering\arraybackslash}X|>{\centering\arraybackslash}X|>{\centering\arraybackslash}X|>{\centering\arraybackslash}X|}
\hline
Collection & Thời gian chạy & Tổng số request & Request/s & Thời gian phản hồi trung bình (ms) \\
\hline
User & 5 phút & 14.799 & 48.02 & 378 \\
\hline
Staff & 5 phút & 12.113 & 39.2 & 333 \\
\hline
Voucher & 5 phút & 7.442 & 24.19 & 341 \\
\hline
Review & 5 phút & 10.992 & 37.72 & 361 \\
\hline
Product & 5 phút & 15.235 & 53.25 & 427 \\
\hline
Order & 5 phút & 8.274 & 31.2 & 391 \\
\hline
Category & 5 phút & 12.643 & 35.19 & 301 \\
\hline
\end{tabularx}
\caption{Bảng số liệu hiệu năng Collection API}
\end{table}


